\documentclass{article}
\usepackage[margin=1in]{geometry}
\usepackage{amsmath, amssymb, amsthm}
\usepackage{enumitem}

%Formatting and Spacing
\usepackage{setspace}
\setenumerate[1]{label = (\alph*), parsep = 1pt, topsep = 5pt}
\setenumerate[2]{label = \roman*.}
\setlength\parindent{0pt}
\linespread{1.10}

\newlist{Cases}{enumerate}{3}
\setlist[Cases]{leftmargin = .25in, label = {Case \arabic*.}, topsep = 0.01in, itemsep = 0.04in, itemindent = .5in, parsep = 0in}

%Custom Commands
\newcommand{\hmwkTitle}{Exam 1 Extra Credit}
\newcommand{\hmwkDueDate}{October 13, 2021}
\newcommand{\hmwkClass}{Honors Discrete Mathematics}
\newcommand{\hmwkClassInstructor}{Professor Gerandy Brito}
\newcommand{\hmwkAuthorName}{Sarthak Mohanty}

%Fancy Headers and Footers
\usepackage{fancyhdr}
\usepackage{extramarks}
\pagestyle{fancy}
\lhead{\hmwkAuthorName}
\chead{\hmwkClass \ (\hmwkClassInstructor)}
\rhead{\hmwkTitle}
\cfoot{\thepage}
\renewcommand\headrulewidth{0.4pt}
\renewcommand\footrulewidth{0.4pt}

\title{
    \vspace{2in}
    \textbf{\hmwkClass:\ \hmwkTitle}\\
    \normalsize\vspace{0.1in}\small{Due\ on\ \hmwkDueDate\ at 11:59pm}\\
    \vspace{0.1in}\large{\textit{\hmwkClassInstructor}}
    \vspace{3in}
    \author{\textbf{\hmwkAuthorName} \\ \\ \small No Collaborators, No Outside Sources}
    \date{}
}

\begin{document}

\maketitle
\pagebreak

\subsection*{Question 1}
    \textit{(Hilbert's Grand Hotel)} The Grand Hotel has countable many rooms, numbered 1, 2, 3, \dots, and it is fully occupied.
    \begin{enumerate}
        \item Suppose that $k$ guests arrive at the hotel. Show how you can accommodate all new guests without removing any of the current guests.
        \item Suppose that a countable infinite number of guests arrive. Show how you can accommodate all new guests without removing any of the current guests.
        \item Suppose Hilbert (the owner of the Grand Hotel) expands the property to a second building, also with infinite rooms numbered $1, 2, 3, \dots$. Show how you can spread the current guests to fill every room in both buildings.
    \end{enumerate}

\subsection*{Solution}
    \begin{enumerate}
        \item Shift all existing guests $k$ rooms over, and place the $k$ guests in the rooms $1, \dots k$.
        \item Move all existing guests in room $n$ to room $2n + 1$. This frees all rooms that can be expressed as $2n$ for the countably infinite number of guests to occupy.
        \item Take every guest with an odd room number $2k - 1$, where $k \in \mathbb{N}$ and move them into room number $k$ in the second building.
        
        \vspace{1.5mm}
        Now take the remaining guests in the original building. Since their room numbers are even, we can represent their room number as $m = 2l$, where $l \in \mathbb{N}$. Move these guests into room number $l$ in the first building.
    \end{enumerate}

\subsection*{Exercise 2}
    Prove that any positive integer can be written as the sum of two or more distinct Fibonacci numbers.

\subsection*{Solution}
    \textsc{Base Case}:  $P(1)$ is true, since $0$ and $1$ are both Fibonacci numbers and $0 + 1 = 1$. \\
    \textsc{Inductive Step}: Now let $n \in \mathbb{N}$ such that $P(1), P(2), \dots P(n)$ are all true. We wish to prove that $P(n + 1)$ is true. Let $f_{n}$ be the largest Fibonacci number such that $n + 1 = f_{n} + r_{n}$, where $0 \le r_{n} < n + 1$.
    \begin{Cases}
        \item Suppose $n + 1 = f_{n}$. Then $n + 1 = f_{n} + 0$, and thus $n + 1$ can be written as the sum of two distinct Fibonacci numbers.
        \item Suppose $n + 1 \ne f_{n}$. Then $P(r_{n})$ is true by the inductive hypothesis, so $n + 1$ can be written as the sum of two or more Fibonacci numbers.
    \end{Cases}
    Since $f_{n} > r_{n}$, we have that $n + 1$ is distinct. Hence $P(n + 1)$ is true. \\
    \textsc{Conclusion}: Therefore by strong induction, for all $n \in \mathbb{N}$, $P(n)$ is true.

\subsection*{Exercise 3}
    Let $n$ be a positive integer and $[n] = 1, \dots, n$. Show that $$\sum_{S \in P([n]) \, \mid \, S \ne \emptyset} \: \prod_{i \in S} \frac{1}{i} = n.$$ Here, $P([n])$ is the power set of $[n]$ and summation is over all the non-empty subsets of $[n]$.

\subsection*{Solution}
    Let $T(n)$ be the sentence $$\sum_{S \in P([n]) \, \mid \, S \ne \emptyset} \: \prod_{i \in S} \frac{1}{i} = n.$$ 
    \textsc{Base Case}:  $T(1)$ is true, since the only non-empty subset of $[n]$ is $\{1\}$ and $\frac{1}{1}$ = 1. \\
    \textsc{Inductive Step}: Now let $n \in \mathbb{N}$ such that $T(n)$ is true. We wish to show that $T(n + 1)$ is true, in other words, we wish to show that $$\sum_{S \in P([n + 1]) \, \mid \, S \ne \emptyset} \: \prod_{i \in S} \frac{1}{i} = n + 1 \quad\text{ or }\quad \left(\sum_{S \in P([n + 1]) \setminus P([n])} \: \prod_{i \in S} \frac{1}{i}\right) + \left(\sum_{S \in P([n]) \, \mid \, S \ne \emptyset} \: \prod_{i \in S} \frac{1}{i}\right) = n + 1.$$ Now from the inductive hypothesis, we reduce this problem to $$\sum_{S \in P([n + 1]) \setminus P([n])} \: \prod_{i \in S} \frac{1}{i} = 1.$$ The set $P([n + 1]) \setminus P([n])$ can be alternatively defined: take the power set $P([n])$ and add the element $n + 1$ to every set in this power set. Hence 
    \begin{align*}
        \sum_{S \in P([n + 1]) \setminus P([n])} \: \prod_{i \in S} \frac{1}{i} &= \frac{1}{n + 1}\sum_{S \in P([n])} \: \prod_{i \in S} \frac{1}{i} \\
        &= \frac{1}{n + 1}\left(\sum_{S \in P([n]) \, \mid \, S \ne \emptyset} \: \prod_{i \in S} \frac{1}{i}\right) + \frac{1}{n + 1}.
    \end{align*}
    Once again from the inductive hypothesis, we know the above expression is equivalent to the value $\frac{n}{n + 1} + \frac{1}{n + 1}= \frac{n + 1}{n + 1} = 1$. Hence $T(n + 1)$ is true as well. \\
    \textsc{Conclusion}: Therefore by induction, for all $n \in \mathbb{N}$, $T(n)$ is true.

\end{document}